% !TEX program = pdflatex
% SIG Sols Togo — Rapport fonctionnel complet
% Compilation : cd docs/rapport && ./compile-rapport.sh
\documentclass[11pt,a4paper]{report}

\usepackage[utf8]{inputenc}
\usepackage[T1]{fontenc}
\usepackage[french]{babel}
\usepackage[margin=2.2cm,headheight=14pt]{geometry}
\usepackage{xcolor}
\usepackage{graphicx}
\usepackage{booktabs}
\usepackage{longtable}
\usepackage{tabularx}
\usepackage{array}
\usepackage{enumitem}
\usepackage{hyperref}
\usepackage{fancyhdr}
\usepackage{titlesec}
\usepackage{tikz}
\usepackage{tcolorbox}
\tcbuselibrary{skins,breakable}

% --- Palette DUSOL ---
\definecolor{DusolVert}{HTML}{0D2818}
\definecolor{DusolVertClair}{HTML}{1A3D2E}
\definecolor{DusolOr}{HTML}{C9A962}
\definecolor{DusolAccent}{HTML}{4ADE80}
\definecolor{DusolFond}{HTML}{F4F7F5}
\definecolor{DusolBleu}{HTML}{1D4ED8}

\hypersetup{
  colorlinks=true,
  linkcolor=DusolVert,
  urlcolor=DusolBleu,
  citecolor=DusolOr,
  pdftitle={SIG Sols Togo — Rapport fonctionnel},
  pdfauthor={DUSOL / DISIA — Maxime Dzidula KELI},
}

\pagestyle{fancy}
\fancyhf{}
\fancyhead[L]{\small\textcolor{DusolVert}{\textbf{SIG Sols Togo}}}
\fancyhead[R]{\small\textcolor{DusolOr}{DUSOL / DISIA}}
\fancyfoot[C]{\thepage}
\renewcommand{\headrulewidth}{0.4pt}
\renewcommand{\footrulewidth}{0.4pt}

\titleformat{\chapter}[display]
  {\normalfont\huge\bfseries\color{DusolVert}}
  {\filleft\Large\color{DusolOr}\chaptertitlename\ \thechapter}
  {0.5ex}
  {\titlerule\vspace{0.5ex}\filright}
  [\vspace{0.5ex}]

\newtcolorbox{infobox}[1]{
  enhanced,
  breakable,
  colback=DusolFond,
  colframe=DusolVertClair,
  fonttitle=\bfseries\color{white},
  coltitle=white,
  attach boxed title to top left={yshift=-2mm,xshift=5mm},
  boxed title style={colback=DusolVert,sharp corners},
  title=#1,
  left=8pt,right=8pt,top=8pt,bottom=8pt,
}

\newtcolorbox{apibox}[1]{
  enhanced,
  breakable,
  colback=white,
  colframe=DusolOr,
  fonttitle=\bfseries,
  title=#1,
  left=6pt,right=6pt,
}

\newcommand{\endpoint}[2]{\texttt{#1} \hfill \textit{#2}\\}

\begin{document}

% ===== PAGE DE GARDE =====
\begin{titlepage}
  \begin{tikzpicture}[remember picture,overlay]
    \fill[DusolVert] (current page.north west) rectangle ([yshift=-8cm]current page.north east);
    \fill[DusolFond] ([yshift=-8cm]current page.north west) rectangle (current page.south east);
    \node[anchor=west,text=DusolOr,font=\Huge\bfseries] at ([xshift=2cm,yshift=-3cm]current page.north west) {SIG Sols Togo};
    \node[anchor=west,text=white,font=\Large] at ([xshift=2cm,yshift=-4.5cm]current page.north west) {Système d'Information Géographique des Sols};
    \node[anchor=west,text=DusolAccent,font=\normalsize] at ([xshift=2cm,yshift=-5.8cm]current page.north west) {Région Maritime · République du Togo};
  \end{tikzpicture}
  \vspace*{9cm}
  {\LARGE\bfseries\color{DusolVert} Rapport fonctionnel détaillé\par}
  \vspace{0.8cm}
  {\large Plateforme web PWA · Application mobile Flutter · API Django/PostGIS\par}
  \vspace{1.5cm}
  \begin{infobox}{Projet et porteur}
    \textbf{DUSOL / DISIA} — Direction des Systèmes d'Information et de l'Analyse\\
    \textbf{Développement :} Maxime Dzidula KELI · +33 754830039\\
    \textbf{Base de données :} PostGIS \texttt{sig\_sols} (partagée web + mobile)\\
    \textbf{URL pilote :} \url{http://localhost:8081}
  \end{infobox}
  \vfill
  {\small\color{gray} Document généré pour documentation technique et présentation institutionnelle.\\
  Juin 2026 — Version 1.0}
\end{titlepage}

\tableofcontents
\listoffigures
\clearpage

% ===== RÉSUMÉ EXÉCUTIF =====
\chapter{Résumé exécutif}

\begin{infobox}{Vision du projet}
SIG Sols Togo est une plateforme géospatiale pédagogique et opérationnelle dédiée à la
\textbf{cartographie, l'analyse et la sensibilisation} autour des sols de la Région Maritime du Togo.
Elle intègre des données terrain (pH, humidité, type de sol), des services spatiaux avancés
(NASA Earthdata, Sentinel Hub, OpenWeather), de l'intelligence artificielle (Gemini, ML fertilité),
ainsi qu'un volet éducatif (quiz, fiches PDF, vidéos, communauté).
\end{infobox}

\section{Objectifs principaux}
\begin{itemize}[leftmargin=*]
  \item Centraliser les \textbf{points sol} géoréférencés dans PostGIS avec validation agent.
  \item Offrir une \textbf{carte interactive} multi-fonds (OSM, satellite, topographique) et multi-couches (NDVI NASA/Sentinel).
  \item Analyser des \textbf{parcelles administratives} (cantons, zones dégradées) ou des polygones dessinés.
  \item Sensibiliser via \textbf{quiz certifiants}, fiches pédagogiques et contenus vidéo.
  \item Permettre le travail \textbf{terrain hors ligne} (agents) avec synchronisation automatique.
  \item Assurer la \textbf{parité web/mobile} sur les mêmes API REST et la même base PostGIS.
\end{itemize}

\section{Schéma d'architecture global}

\begin{figure}[h]
\centering
\begin{tikzpicture}[
  node distance=1.1cm,
  box/.style={rectangle,rounded corners,draw=DusolVert,fill=DusolFond,minimum width=3.2cm,minimum height=0.9cm,align=center,font=\small},
  srv/.style={rectangle,rounded corners,draw=DusolOr,fill=white,minimum width=3.5cm,minimum height=0.9cm,align=center,font=\small},
]
  \node[box] (web) {Site web PWA\\Leaflet + JS};
  \node[box,right=of web] (mob) {App mobile\\Flutter};
  \node[srv,below=1.4cm of web,xshift=1.8cm] (api) {API Django REST\\+ Daphne WebSocket};
  \node[srv,below=of api] (db) {PostgreSQL/PostGIS};
  \node[srv,left=2.8cm of api] (ext) {APIs externes\\NASA · Sentinel · OW · Gemini};
  \node[srv,right=2.8cm of api] (cel) {Celery + Redis\\tâches planifiées};
  \draw[->,thick,DusolVert] (web) -- (api);
  \draw[->,thick,DusolVert] (mob) -- (api);
  \draw[->,thick,DusolVert] (api) -- (db);
  \draw[->,thick,DusolOr] (api) -- (ext);
  \draw[->,thick,DusolOr] (cel) -- (api);
  \draw[->,thick,DusolOr] (cel) -- (db);
\end{tikzpicture}
\caption{Architecture trois-tiers : clients web et mobile, backend unifié, PostGIS central.}
\end{figure}

% ===== INFRASTRUCTURE =====
\chapter{Infrastructure et déploiement}

\section{Stack technique}
\begin{tabularx}{\textwidth}{@{}lX@{}}
\toprule
\textbf{Composant} & \textbf{Rôle} \\
\midrule
PostgreSQL 15 + PostGIS 3.4 & Base spatiale \texttt{sig\_sols} \\
Redis 7 & Cache, Celery, Channels (WebSocket) \\
Django 4 + GeoDjango + DRF & API REST, admin, ORM spatial \\
Daphne (ASGI) & HTTP + WebSocket \texttt{/ws/live/} \\
Celery + Beat & Ingestion NASA, alertes sécheresse, ML, purge GPS \\
NGINX & Reverse proxy :8081, static, WebSocket upgrade \\
Docker Compose & Orchestration multi-conteneurs \\
\bottomrule
\end{tabularx}

\section{Services Docker}
\begin{itemize}[leftmargin=*]
  \item \texttt{db} — PostGIS persistant
  \item \texttt{web} — Daphne sur port 8000 interne
  \item \texttt{nginx} — exposition publique port 8081
  \item \texttt{redis}, \texttt{celery}, \texttt{celery-beat}
\end{itemize}

\section{Comptes de démonstration}
\begin{tabular}{lll}
\toprule
Utilisateur & Mot de passe & Rôle \\
\midrule
\texttt{admin} & \texttt{admin123} & Administrateur \\
\texttt{agent1} & \texttt{agent123} & Agent terrain \\
\texttt{public1} & \texttt{public123} & Utilisateur public \\
\bottomrule
\end{tabular}

% ===== CARTE ET SPATIAL =====
\chapter{Cartographie interactive}

\section{Fonds de carte (3 topographies / basemaps)}
Le site web et l'application mobile proposent \textbf{trois fonds de carte} interchangeables :

\begin{infobox}{Les 3 types de fonds cartographiques}
\begin{enumerate}
  \item \textbf{OpenStreetMap (Plan)} — rue, limites, toponymie standard.
  \item \textbf{Satellite (Esri World Imagery)} — imagerie aérienne haute résolution.
  \item \textbf{Topographique (OpenTopoMap)} — courbes de niveau, relief, sentiers.
\end{enumerate}
\end{infobox}

\section{Couches et outils carte (web)}
\begin{itemize}[leftmargin=*]
  \item Points sol colorés par pH (rouge/orange/vert/bleu)
  \item Clustering de marqueurs, heatmap pH
  \item Tuiles NDVI NASA et Sentinel Hub (opacité réglable)
  \item Badge météo OpenWeather sur la carte
  \item Filtres : pH min/max, type de sol, validés, bbox visible
  \item Export GeoJSON / CSV ; import GeoJSON/CSV (agents)
  \item Ajout point au clic, file d'attente offline
  \item Proximité GPS, comparaison 2 points, trajectoire 24h
  \item Partage position live + WebSocket \texttt{/ws/live/}
  \item Alertes sécheresse et notifications de proximité
\end{itemize}

\section{Parcelles et zones administratives}

\begin{infobox}{Sélection de parcelles}
Les \textbf{parcelles} correspondent aux \textbf{zones administratives} PostGIS :
\begin{itemize}
  \item \textbf{Cantons} (\texttt{zone\_type=canton}) — unités pilotes Région Maritime
  \item \textbf{Zones dégradées} (\texttt{zone\_type=degraded}) — surfaces à risque
  \item Filtres : Toutes / Cantons / Dégradées
  \item Sélection : liste déroulante, clic sur carte, ou dessin polygone (web)
\end{itemize}
\end{infobox}

\subsection{Analyse parcelle live}
L'endpoint \texttt{POST /api/v1/spatial/parcel/live/} agrège :
\begin{itemize}
  \item Statistiques sols dans la zone (pH moyen, points validés)
  \item \textbf{NASA} : NDVI MOD13Q1, humidité SMAP
  \item \textbf{Sentinel Hub} : NDVI moyen Sentinel-2
  \item \textbf{OpenWeather} : météo actuelle et prévisions
  \item \textbf{ML} : prédiction fertilité (RandomForest/XGBoost)
  \item Indice de santé, vulnérabilité, recommandations agronomiques
\end{itemize}

% ===== API BACKEND =====
\chapter{API Backend — Inventaire des endpoints}

\section{Authentification \texttt{/api/v1/auth/}}
\begin{apibox}{Comptes, profil, social, GPS}
\endpoint{POST /auth/register/}{Inscription étendue}
\endpoint{POST /auth/token/}{JWT access + refresh + profil}
\endpoint{POST /auth/token/refresh/}{Renouvellement token}
\endpoint{GET/PATCH /auth/profile/}{Profil utilisateur}
\endpoint{POST /auth/profile/photo/}{Photo de profil}
\endpoint{POST /auth/password/change/}{Changement mot de passe}
\endpoint{GET /auth/feed/}{Fil abonnements}
\endpoint{GET/POST/DELETE /auth/favorites/}{Favoris (fiches, vidéos, points)}
\endpoint{GET/POST /auth/messages/}{Messages directs}
\endpoint{POST /auth/location/}{Publier position GPS}
\endpoint{GET /auth/locations/live/}{Agents en ligne}
\endpoint{GET /auth/trajectory/}{Historique GPS 24h}
\end{apibox}

\section{Points sol \texttt{/api/v1/points/}}
\begin{apibox}{CRUD spatial et validation}
\endpoint{GET/POST /points/}{Liste et création GeoJSON}
\endpoint{GET/PATCH/DELETE /points/<id>/}{Détail point}
\endpoint{POST /points/<id>/validate\_point/}{Valider/rejeter (admin)}
\endpoint{GET /points/geojson/}{Export GeoJSON}
\endpoint{GET /points/export-csv/}{Export CSV}
\endpoint{POST /points/import\_data/}{Import fichier (agent)}
\endpoint{GET /points/compare/}{Comparer deux points}
\endpoint{GET /heatmap/}{Données heatmap pH}
\endpoint{GET /dashboard/stats/}{KPIs tableau de bord}
\endpoint{GET /validation/pending/}{File validation admin}
\end{apibox}

\section{Spatial \texttt{/api/v1/spatial/}}
\begin{apibox}{GIS, parcelles, vulnérabilité}
\endpoint{GET /spatial/proximity/}{Points dans un rayon}
\endpoint{POST /spatial/intersection/}{Intersection géométrie/zones}
\endpoint{POST /spatial/buffer/}{Tampon}
\endpoint{POST /spatial/area/}{Surface polygone}
\endpoint{GET /spatial/vulnerability/}{Zonage vulnérabilité}
\endpoint{GET /spatial/ndvi-timeseries/<id>/}{Série NDVI NASA}
\endpoint{GET /spatial/statistics/}{Stats par canton}
\endpoint{GET /spatial/smap-correlation/}{Corrélation SMAP terrain}
\endpoint{POST /spatial/parcel/analyze/}{Analyse complète}
\endpoint{GET/POST /spatial/parcel/live/}{Analyse temps réel}
\endpoint{GET /spatial/parcel/zones/}{Liste zones}
\endpoint{GET /spatial/parcel/zones/geojson/}{GeoJSON cantons/dégradées}
\end{apibox}

\section{APIs externes intégrées}
\begin{tabularx}{\textwidth}{@{}lXl@{}}
\toprule
\textbf{Service} & \textbf{Fonction} & \textbf{Préfixe API} \\
\midrule
NASA Earthdata / STAC & Catalogue, ingestion, tuiles NDVI & \texttt{/nasa/} \\
Sentinel Hub & OAuth, NDVI, true color, tuiles & \texttt{/sentinel/} \\
OpenWeatherMap & Actuel, prévisions, statut & \texttt{/weather/} \\
Google Gemini & Chatbot pédagogique & \texttt{/assistant/} \\
ML local & Fertilité RF/XGBoost & \texttt{/ml/} \\
\bottomrule
\end{tabularx}

\section{Éducation \texttt{/api/v1/education/}}
\begin{itemize}[leftmargin=*]
  \item Quiz : démarrer, répondre, terminer, classement, badges, défi hebdo
  \item Certificat PDF avec token de vérification
  \item 15+ fiches pédagogiques (PDF ReportLab)
  \item Parcours d'apprentissage adaptatif
\end{itemize}

\section{Vidéos \texttt{/api/v1/videos/}}
\begin{itemize}[leftmargin=*]
  \item Vidéos et shorts (upload, likes, commentaires, vues)
  \item Stories éphémères 24h
  \item Workflow modération : pending → published/rejected
  \item Modération commentaires (masquer, hint IA)
\end{itemize}

\section{Plateforme \texttt{/api/v1/platform/}}
\begin{itemize}[leftmargin=*]
  \item Notifications, recherche globale, dashboard personnel
  \item Admin : KPIs, analytics, exports CSV, rapports ministère/zone
  \item Alertes sécheresse, audit, reset mot de passe
  \item Suivi d'activité client \texttt{POST /platform/activity/}
\end{itemize}

% ===== WEB =====
\chapter{Application web (PWA)}

\section{Écrans et navigation}
\begin{tabularx}{\textwidth}{@{}lX@{}}
\toprule
\textbf{Écran} & \textbf{Contenu principal} \\
\midrule
Carte & Filtres, parcelles, outils, couches NASA/Sentinel/météo \\
Tableau de bord & KPIs, graphiques Chart.js, corrélation SMAP \\
Quiz & QCM chronométrés, certificats, badges, classement \\
Fiches & PDF intégrés, favoris \\
Vidéos / Shorts & Upload, modération, stories \\
Communauté & Fil, favoris, messages, profils publics \\
Mon espace & Stats personnelles, parcours \\
Admin & Validation, analytics, exports, ML/NASA \\
Aide & FAQ, visite guidée, raccourcis clavier \\
\bottomrule
\end{tabularx}

\section{Fonctionnalités transverses web}
\begin{itemize}[leftmargin=*]
  \item \textbf{Offline} : file \texttt{sig\_sols\_offline\_queue} (localStorage)
  \item \textbf{WebSocket} : positions live groupe \texttt{sig\_sols\_live}
  \item \textbf{i18n} : FR, EN, Éwé, Kabyle
  \item \textbf{Thème} : clair / sombre
  \item \textbf{PWA} : manifest, service worker cache statique
  \item \textbf{Assistant Gemini} : FAB flottant contextuel
\end{itemize}

% ===== MOBILE =====
\chapter{Application mobile Flutter}

\section{Parité avec le web}
L'application mobile (\texttt{mobile/}) consomme les \textbf{mêmes endpoints} via \texttt{SigApi} et \texttt{ApiClient} JWT.

\section{Écrans mobiles}
\begin{tabularx}{\textwidth}{@{}ll@{}}
\toprule
\textbf{Route} & \textbf{Fonctionnalités} \\
\midrule
\texttt{/} Carte & 3 fonds, parcelles, NDVI, outils, offline, live GPS \\
\texttt{/dashboard} & KPIs, cartes APIs externes \\
\texttt{/quiz} & Quiz complet, badges, défi hebdo \\
\texttt{/sheets} & Fiches PDF + favoris \\
\texttt{/videos}, \texttt{/shorts} & Upload, likes, commentaires \\
\texttt{/community} & Fil, favoris, messages \\
\texttt{/assistant} & Chat Gemini \\
\texttt{/profile} & Édition profil, photo, trajectoire \\
\texttt{/parcel} & Analyse zone canton + GPS \\
\texttt{/search} & Recherche globale \\
\texttt{/notifications} & Badge + liste \\
\texttt{/my-dashboard} & Espace personnel \\
\texttt{/admin} & KPIs, validation, exports (admin) \\
\texttt{/help} & Aide \\
\bottomrule
\end{tabularx}

\section{Spécificités mobiles}
\begin{itemize}[leftmargin=*]
  \item Synchronisation offline points sol (SharedPreferences)
  \item WebSocket JWT \texttt{?token=} + polling GPS fallback
  \item i18n FR/EN, thème clair/sombre
  \item Activity tracker vers \texttt{/platform/activity/}
\end{itemize}

% ===== ADMIN =====
\chapter{Administration et modération}

\section{Panneau admin web}
\begin{itemize}[leftmargin=*]
  \item KPIs temps réel : utilisateurs, points, agents live, alertes
  \item Validation points sol en attente
  \item Modération vidéos et commentaires
  \item Analytics : heatmap activité, top utilisateurs, timeline
  \item Exports CSV utilisateurs et activité 30 jours
  \item Rapport ministère HTML, rapport par zone
  \item Réentraînement ML, ingestion NASA manuelle
\end{itemize}

% ===== TÂCHES PLANIFIÉES =====
\chapter{Tâches planifiées (Celery)}

\begin{tabularx}{\textwidth}{@{}lX@{}}
\toprule
\textbf{Tâche} & \textbf{Fréquence} \\
\midrule
Ingestion NASA Earthdata & Hebdomadaire (lundi 02h00) \\
Vérification réentraînement ML & Mensuelle \\
Alertes sécheresse & Horaire \\
Purge historique GPS & Quotidienne \\
Rafraîchissement classement quiz & Configurable \\
\bottomrule
\end{tabularx}

% ===== SYNTHÈSE =====
\chapter{Synthèse des modules}

\begin{figure}[h]
\centering
\begin{tikzpicture}[
  module/.style={rectangle,rounded corners,draw=DusolVert,fill=DusolFond,minimum width=2.6cm,minimum height=0.8cm,align=center,font=\footnotesize},
]
  \node[module] (m1) at (0,0) {Carte\\Spatial};
  \node[module] (m2) at (3,0) {Parcelles\\Zones};
  \node[module] (m3) at (6,0) {NASA\\Sentinel};
  \node[module] (m4) at (9,0) {Météo\\ML};
  \node[module] (m5) at (0,-1.5) {Quiz\\Fiches};
  \node[module] (m6) at (3,-1.5) {Vidéos\\Stories};
  \node[module] (m7) at (6,-1.5) {Communauté\\Messages};
  \node[module] (m8) at (9,-1.5) {Admin\\Audit};
  \node[module,fill=DusolOr!30] (c) at (4.5,-3) {PostGIS sig\_sols};
  \foreach \n in {m1,m2,m3,m4,m5,m6,m7,m8} \draw[->,thick] (\n) -- (c);
\end{tikzpicture}
\caption{Modules fonctionnels convergent vers la base PostGIS unique.}
\end{figure}

\chapter*{Conclusion}
\addcontentsline{toc}{chapter}{Conclusion}

SIG Sols Togo constitue une plateforme complète couvrant la chaîne de valeur
\textbf{collecte terrain → analyse spatiale multi-sources → pédagogie → gouvernance admin},
avec une architecture moderne (Docker, PostGIS, JWT, WebSocket, Celery) et une
\textbf{double interface web/mobile} partageant la même source de vérité.

\vspace{1cm}
\begin{center}
\textcolor{DusolVert}{\textbf{DUSOL / DISIA — SIG Sols Togo}}\\
\textcolor{DusolOr}{Maxime Dzidula KELI · +33 754830039}
\end{center}

\end{document}
